\section{Introduction}

\lettrine[lines=3,lhang=0.08,loversize=0.15,findent=0.5em,nindent=0em]{F}{irmware} update mechanisms are one of the most security-critical components of modern embedded systems. Routers, secure elements, industrial controllers, vehicles, and consumer IoT devices all rely on firmware to initialize hardware and enforce privileged security decisions. If an attacker can install unauthorized firmware, the entire higher-level security architecture can be bypassed. For that reason, secure boot and update systems typically require a digital signature before code is executed or accepted for installation.

This trust model is under long-term pressure from quantum computing. NIST finalized the first post-quantum cryptography standards on August 13, 2024, including FIPS 204 for ML-DSA and FIPS 205 for SLH-DSA, and encouraged organizations to begin migrating to the new standards as soon as possible \cite{nist_pqc_release,fips204,fips205}. The migration question is especially urgent for firmware because hardware development and deployment cycles are long. As observed in the OpenTitan secure-boot context, a boot ROM designed today may remain fielded for ten to twenty years, and some early verification code can never be changed after manufacture \cite{opentitan_secure_boot,zerorisc_opentitan}.

However, firmware signing is not merely a benchmark problem about key size and verification time. A realistic update system must also authenticate metadata, prevent rollback to older vulnerable images, bind images to the intended device class, and avoid migration policies that allow attackers to force a fallback to weaker signatures. A package that carries a valid post-quantum signature but leaves the algorithm identifier or minimum bootloader version unsigned can still fail to provide an adequate security boundary.

This paper therefore studies post-quantum firmware signing as a systems problem rather than as a narrow primitive replacement. The focus is on ML-DSA and SLH-DSA because NIST positions ML-DSA as the primary digital signature standard and SLH-DSA as a backup based on a different mathematical approach \cite{nist_pqc_release}. The same migration pressure is visible in industry firmware guidance such as the UEFI Forum's post-quantum update whitepaper \cite{uefi_pqc_whitepaper}. Classical baselines, especially RSA-PSS and ECDSA P-256, are included because they remain dominant in legacy firmware ecosystems and define the practical migration target.

The main contributions of this work are as follows:
\begin{enumerate}
\item A firmware-oriented comparative analysis of RSA-PSS, ECDSA P-256, ML-DSA, and SLH-DSA that emphasizes signature size, key size, conservative security assumptions, and bootloader suitability rather than raw cryptographic novelty.
\item A proposed PQC-ready firmware package design that signs both metadata and firmware content, supports anti-rollback and algorithm agility, and explicitly binds accepted signature algorithms into authenticated policy state.
\item A modular Python prototype that models firmware package metadata and migration policies, generates reproducible evaluation tables, and performs real ML-DSA and SPHINCS+-family signing and verification through local PQC bindings.
\item An evaluation that quantifies signature overhead for firmware images from 64 KB to 4 MB and simulates how five migration policies handle downgrade, rollback, and tampering cases.
\end{enumerate}

The remainder of the paper is organized as follows. Section II reviews firmware update, secure boot, and the relevant signature families. Section III defines the threat model and security requirements. Section IV compares classical and post-quantum signature options. Section V proposes the firmware package structure and verification logic. Section VI describes the modular prototype. Section VII presents the evaluation results. Section VIII discusses deployment trade-offs, and Section IX concludes the paper.
